\documentclass[11pt,a4paper]{beamer}
\usetheme{Madrid}
\usepackage[utf8]{inputenc}
\usepackage{amsmath}
\usepackage{amsfonts}
\usepackage{amssymb}
\usepackage{graphicx}
\usepackage{tikz}
\usepackage{booktabs}
\usepackage{array}
\usepackage{colortbl}
\usepackage{multicol}

\usefonttheme{professionalfonts}

\setbeamerfont{normal text}{size=\tiny}
\setbeamerfont{itemize/enumerate body}{size=\tiny}
\setbeamerfont{itemize/enumerate subbody}{size=\tiny}
\setbeamerfont{block body}{size=\tiny}
\setbeamerfont{block title}{size=\scriptsize}
\setbeamerfont{frametitle}{size=\small}
\setbeamerfont{title}{size=\small}
\setbeamerfont{subtitle}{size=\footnotesize}
\setbeamerfont{author}{size=\tiny}
\setbeamerfont{date}{size=\tiny}

\setlength{\leftmargini}{0.3em}
\setlength{\leftmarginii}{0.3em}
\setlength{\labelsep}{0.1em}
\setlength{\itemsep}{0.02em}
\setlength{\parskip}{0.02em}

\title[CAMS Exam Preparation]{CAMS Exam Preparation}
\subtitle{Money Laundering Risks in DNFBPs}
\author{Lawyers | Accountants | Real Estate | Casinos | Trust Companies}
\date{}

\setbeamercolor{title}{fg=white,bg=blue!70!black}
\setbeamercolor{frametitle}{fg=white,bg=blue!70!black}
\setbeamertemplate{navigation symbols}{}

\begin{document}

% TITLE SLIDE
\begin{frame}
\titlepage
\centering
\tiny 15 Topics | DNFBP ML Risks | How to Answer
\end{frame}

% ============================================================
% SLIDE 1 - DNFBP OVERVIEW AND FATF REQUIREMENTS
% ============================================================
\begin{frame}{1. DNFBP Overview and FATF Requirements}
\tiny
\noindent\textbf{Introduction to the topic:} Designated Non-Financial Businesses and Professions (DNFBPs) are sectors that, while not traditional financial institutions, are vulnerable to money laundering and terrorist financing due to their role in facilitating financial transactions, holding client assets, or providing professional services that can be misused. Under FATF Recommendations 22 and 23, DNFBPs include: casinos (including internet casinos), real estate agents, dealers in precious metals and stones, lawyers, notaries, and other independent legal professionals when they engage in activities such as buying or selling real estate, managing client money or assets, managing bank accounts, or creating legal entities. Accountants and trust and company service providers are also included.

\vspace{0.1cm}
\noindent\textbf{Type of violation or warning signs:} A lawyer or notary accepts large cash payments without verifying the client's identity or source of funds. A real estate agent facilitates a property purchase through a shell company without identifying the beneficial owner. A casino allows high-value gamblers to use cash without proper due diligence or transaction reporting. A trust company creates a complex structure with anonymous beneficiaries and does not identify the source of the assets. A dealer in precious metals accepts large cash payments without filing a Currency Transaction Report or SAR.

\vspace{0.1cm}
\noindent\textbf{Correct answer approach:} The correct answer will state that DNFBPs are subject to AML/CFT obligations including Customer Due Diligence, recordkeeping, and reporting of suspicious transactions. Eliminate answers that suggest DNFBPs are exempt from AML requirements or that professional privilege excuses them from all obligations. Remember: FATF Recommendations 22 and 23 specifically target DNFBPs; they must have AML programs and file SARs.
\end{frame}

% ============================================================
% SLIDE 2 - LAWYERS AND NOTARIES ML RISKS
% ============================================================
\begin{frame}{2. Lawyers and Notaries ML Risks}
\tiny
\noindent\textbf{Introduction to the topic:} Lawyers and notaries are vulnerable to money laundering because they handle large sums of client money, create complex legal structures (trusts, companies, foundations), facilitate real estate transactions, and manage accounts for clients. Criminals use lawyers to create opaque structures that hide beneficial ownership, move funds through client accounts, or facilitate purchases of high-value assets. While attorney-client privilege is important, it does not exempt lawyers from AML obligations when they are acting in financial transactions. Many jurisdictions require lawyers to report suspicious transactions unless the information is subject to legal privilege.

\vspace{0.1cm}
\noindent\textbf{Type of violation or warning signs:} A lawyer accepts a large cash payment from a client to purchase real estate but does not verify the client's identity or source of funds. A lawyer creates a complex trust structure with anonymous beneficiaries and no clear legitimate purpose. A lawyer's client account receives funds from multiple unrelated sources and is quickly disbursed to third parties. A lawyer refuses to disclose the beneficial owner of a company they are setting up. A lawyer advises a client on how to structure transactions to avoid reporting thresholds.

\vspace{0.1cm}
\noindent\textbf{Correct answer approach:} The correct answer will state that lawyers must conduct CDD and file SARs when engaging in financial transactions, subject to legal privilege where applicable. Eliminate answers that suggest legal privilege exempts lawyers from all AML obligations or that lawyers never need to verify source of funds. Remember: when a lawyer acts as a financial intermediary, AML obligations apply.
\end{frame}

% ============================================================
% SLIDE 3 - REAL ESTATE AGENTS ML RISKS
% ============================================================
\begin{frame}{3. Real Estate Agents ML Risks}
\tiny
\noindent\textbf{Introduction to the topic:} Real estate is a classic vehicle for money laundering because it is high-value, relatively illiquid, and can be purchased through complex structures. Real estate agents are often the first point of contact in a property transaction and must conduct due diligence on their clients. Red flags in real estate include cash purchases, purchases through shell companies or trusts where the beneficial owner is not disclosed, purchases at prices significantly above or below market value, and rapid resale of properties at a loss. Under FATF, real estate agents must identify their clients and verify the source of funds for transactions.

\vspace{0.1cm}
\noindent\textbf{Type of violation or warning signs:} A client purchases a \$5 million property in cash with no mortgage. The client uses a shell company to purchase the property and refuses to disclose the beneficial owner. The client buys a property for \$10 million and sells it six months later for \$7 million — a 30\% loss with no legitimate explanation. The client asks the real estate agent to split the payment into multiple smaller amounts to avoid reporting thresholds. The client is a foreign PEP but the agent does not conduct EDD.

\vspace{0.1cm}
\noindent\textbf{Correct answer approach:} The correct answer will state that real estate agents must conduct CDD on buyers and sellers, identify beneficial owners of any legal entity used, and verify source of funds. Eliminate answers that suggest real estate agents have no AML obligations or that cash purchases are always acceptable. Remember: rapid resale at a loss is a classic money laundering indicator — the loss is the laundering cost.
\end{frame}

% ============================================================
% SLIDE 4 - CASINOS AND GAMBLING ML RISKS
% ============================================================
\begin{frame}{4. Casinos and Gambling ML Risks}
\tiny
\noindent\textbf{Introduction to the topic:} Casinos and gambling establishments are high-risk for money laundering because they handle large amounts of cash, offer anonymity to high-value clients, and have mechanisms to convert cash into chips and then back into cash or checks. Criminals use casinos to layer illicit funds through gambling activities — they may buy chips with dirty money, gamble a small amount, and then cash out the remaining chips as clean winnings. Internet casinos present additional risks because they operate across borders and may be subject to weak AML oversight in some jurisdictions.

\vspace{0.1cm}
\noindent\textbf{Type of violation or warning signs:} A customer buys chips with large amounts of cash, gambles a small portion, and cashes out the remainder as "winnings." A customer purchases chips just under the reporting threshold multiple times (structuring). A customer uses multiple accounts or loyalty cards to fragment transactions. A customer from a high-risk jurisdiction engages in high-value gambling with no apparent source of income. A casino does not file SARs on suspicious patrons. An internet casino accepts payments through unregulated cryptocurrency without verification.

\vspace{0.1cm}
\noindent\textbf{Correct answer approach:} The correct answer will state that casinos must conduct CDD on customers, monitor transactions, file CTRs for cash transactions over the applicable threshold, and file SARs for suspicious activity. Eliminate answers that suggest gambling is an exception to AML rules or that casinos have no obligation to verify source of funds. Remember: the "gambling" is often just a cover for cleaning money — the criminal accepts losses as the laundering cost.
\end{frame}

% ============================================================
% SLIDE 5 - DEALERS IN PRECIOUS METALS AND STONES (DPMs)
% ============================================================
\begin{frame}{5. Dealers in Precious Metals and Stones (DPMs)}
\tiny
\noindent\textbf{Introduction to the topic:} Dealers in precious metals and stones (gold, silver, diamonds, gemstones) are vulnerable to money laundering because precious metals and stones are high-value, portable, easily concealable, and can be traded across borders with minimal documentation. Criminals use precious metals to convert cash into a physical asset that can be sold later for clean funds. Gold and diamonds can be easily smuggled and are difficult to trace. Under FATF, DPMs must conduct customer due diligence on transactions above a certain threshold and report suspicious activity.

\vspace{0.1cm}
\noindent\textbf{Type of violation or warning signs:} A customer purchases gold bars with large amounts of cash without providing identification. A customer purchases diamonds with funds from a high-risk jurisdiction. A customer makes multiple purchases just under the reporting threshold (structuring). A customer sells gold or diamonds and asks to be paid in a different currency or sent to an offshore account. A customer cannot explain the source of the funds or the origin of the precious metals. The DPM does not have an AML program or has never filed a SAR.

\vspace{0.1cm}
\noindent\textbf{Correct answer approach:} The correct answer will state that DPMs must conduct CDD, file CTRs for cash transactions over the reporting threshold, file SARs for suspicious activity, and maintain records. Eliminate answers that suggest DPMs are exempt from AML obligations or that precious metals transactions are always legitimate. Remember: precious metals are a classic money laundering vehicle — the anonymity and portability make them attractive to criminals.
\end{frame}

% ============================================================
% SLIDE 6 - TRUST AND COMPANY SERVICE PROVIDERS (TCSPs)
% ============================================================
\begin{frame}{6. Trust and Company Service Providers (TCSPs)}
\tiny
\noindent\textbf{Introduction to the topic:} Trust and company service providers (TCSPs) create and manage legal entities — trusts, foundations, corporations, and partnerships — that can be used to hide beneficial ownership. TCSPs are at the front line of preventing the misuse of legal entities for money laundering. Under FATF Recommendation 24, countries must ensure that adequate and accurate beneficial ownership information is held by or on behalf of companies. TCSPs must identify the settlor, trustees, and beneficiaries of trusts, and the shareholders and directors of companies they form.

\vspace{0.1cm}
\noindent\textbf{Type of violation or warning signs:} A TCSP creates a trust with anonymous beneficiaries and refuses to disclose them to the bank. A TCSP forms a company for a client from a high-risk jurisdiction with no legitimate business purpose. A TCSP uses nominee directors and shareholders to hide the true beneficial owner. A TCSP creates multiple entities for the same client that appear to have no economic purpose. A TCSP does not maintain records of the beneficial owners of the entities it creates. A TCSP has no AML program and has never filed a SAR.

\vspace{0.1cm}
\noindent\textbf{Correct answer approach:} The correct answer will state that TCSPs must identify beneficial owners of all entities they create and maintain accurate records. Eliminate answers that suggest TCSPs have no obligation to know the beneficial owners or that nominee arrangements are acceptable. Remember: TCSPs are key to transparency — without their cooperation, beneficial ownership cannot be identified.
\end{frame}

% ============================================================
% SLIDE 7 - ACCOUNTANTS AND AUDITORS ML RISKS
% ============================================================
\begin{frame}{7. Accountants and Auditors ML Risks}
\tiny
\noindent\textbf{Introduction to the topic:} Accountants and auditors can be misused by criminals to create false financial statements, hide assets, or facilitate tax evasion and money laundering. They may prepare accounts for shell companies, advise on complex structures that obscure beneficial ownership, or help clients move funds across borders. Under FATF, accountants and auditors are subject to AML obligations when they engage in financial transactions on behalf of clients. They must conduct CDD, maintain records, and report suspicious activity.

\vspace{0.1cm}
\noindent\textbf{Type of violation or warning signs:} An accountant prepares financial statements for a shell company with no legitimate business operations. An auditor does not question unusual transactions or unexplained assets in a client's accounts. An accountant advises a client on how to avoid tax reporting or move funds to offshore accounts without disclosure. An accountant accepts large cash payments without verifying the source. An auditor signs off on financial statements that do not reflect the client's true financial position. The accountant has never filed a SAR despite handling high-risk clients.

\vspace{0.1cm}
\noindent\textbf{Correct answer approach:} The correct answer will state that accountants and auditors must conduct CDD, maintain records, and report suspicious transactions when they act as financial intermediaries. Eliminate answers that suggest auditors are exempt from AML obligations because their work is "technical." Remember: accountants who prepare financial statements for suspicious entities may be facilitating money laundering.
\end{frame}

% ============================================================
% SLIDE 8 - LEGAL PROFESSION PRIVILEGE AND AML
% ============================================================
\begin{frame}{8. Legal Professional Privilege and AML Obligations}
\tiny
\noindent\textbf{Introduction to the topic:} Legal professional privilege (also called attorney-client privilege) is a fundamental protection that allows clients to communicate freely with their lawyers. However, privilege does not exempt lawyers from all AML obligations. When a lawyer is acting as a financial intermediary — handling client money, managing accounts, facilitating transactions — AML obligations apply. The privilege does not protect communications that are in furtherance of a crime, nor does it protect information about the source of funds or the identity of the beneficial owner. Many jurisdictions have explicit exceptions to privilege for AML reporting.

\vspace{0.1cm}
\noindent\textbf{Type of violation or warning signs:} A lawyer claims legal privilege to avoid disclosing the source of client funds used for a property purchase. A lawyer refuses to identify a client's beneficial owner on the basis of privilege. A lawyer uses privilege to shield illegal activity. A lawyer does not have procedures to distinguish between privileged and non-privileged information for AML purposes. A lawyer advises a client to structure transactions to avoid reporting thresholds, claiming it is "legitimate tax planning."

\vspace{0.1cm}
\noindent\textbf{Correct answer approach:} The correct answer will state that legal privilege does not exempt lawyers from AML obligations when they are acting as financial intermediaries, and privilege does not apply to communications that further criminal activity. Eliminate answers that suggest privilege completely exempts lawyers from AML obligations. Remember: privilege applies to legal advice, not to financial transactions. Source of funds is not privileged.
\end{frame}

% ============================================================
% SLIDE 9 - DNFBP SAR FILING OBLIGATIONS
% ============================================================
\begin{frame}{9. DNFBP SAR Filing Obligations}
\tiny
\noindent\textbf{Introduction to the topic:} DNFBPs are required to file Suspicious Activity Reports (SARs) when they suspect money laundering or terrorist financing. This applies to all DNFBPs, not just financial institutions. The threshold for suspicion is the same as for banks — reasonable grounds to suspect. DNFBPs must file SARs regardless of the amount of the transaction. Failure to file SARs is a violation that can result in fines and enforcement actions. DNFBPs must have internal procedures for identifying suspicious activity and escalating it to the compliance officer and to the FIU.

\vspace{0.1cm}
\noindent\textbf{Type of violation or warning signs:} A DNFBP has never filed a SAR despite handling high-risk clients and transactions. A lawyer suspects a client is laundering money but does not file a SAR because the client is "important" or "high-value." A real estate agent notices red flags but does not escalate them internally. A casino has no procedure for identifying suspicious activity. A DNFBP files SARs only when regulators request them, not proactively. A DNFBP does not know who its AML compliance officer is or how to file a SAR.

\vspace{0.1cm}
\noindent\textbf{Correct answer approach:} The correct answer will state that DNFBPs must file SARs when they have reasonable grounds to suspect money laundering. Eliminate answers that suggest SAR filing is optional or that DNFBPs are exempt. Remember: the same SAR obligations apply to DNFBPs as to financial institutions. Failure to file is a regulatory violation.
\end{frame}

% ============================================================
% SLIDE 10 - DNFBP CUSTOMER DUE DILIGENCE (CDD)
% ============================================================
\begin{frame}{10. DNFBP Customer Due Diligence (CDD)}
\tiny
\noindent\textbf{Introduction to the topic:} DNFBPs must conduct Customer Due Diligence (CDD) when establishing a business relationship, when conducting transactions above a certain threshold, or when there are concerns about the client's identity. CDD includes identifying the client, verifying the client's identity using reliable and independent sources, identifying the beneficial owner of any legal entity, and understanding the nature and purpose of the business relationship. Enhanced Due Diligence (EDD) is required for PEPs and other high-risk clients. CDD records must be retained for at least five years.

\vspace{0.1cm}
\noindent\textbf{Type of violation or warning signs:} A DNFBP does not verify the identity of its clients. A real estate agent accepts a purchase through a shell company without identifying the beneficial owner. A lawyer opens a trust account for a client but does not obtain any identification documents. A DNFBP does not conduct EDD on PEPs. A DNFBP does not maintain CDD records or discards them after a short period. A DNFBP relies entirely on a third party to conduct CDD without ensuring that third party has adequate procedures.

\vspace{0.1cm}
\noindent\textbf{Correct answer approach:} The correct answer will state that DNFBPs must conduct CDD on all clients, identify beneficial owners, and maintain records. Eliminate answers that suggest CDD is optional or that DNFBPs can rely entirely on third parties without verification. Remember: CDD is the first line of defense. Without knowing the client, you cannot identify suspicious activity.
\end{frame}

% ============================================================
% SLIDE 11 - DNFBP RECORDKEEPING REQUIREMENTS
% ============================================================
\begin{frame}{11. DNFBP Recordkeeping Requirements}
\tiny
\noindent\textbf{Introduction to the topic:} DNFBPs must maintain records of all transactions and CDD information for at least five years. Records must include identification documents, transaction details, correspondence, and SARs filed. Records must be sufficient to reconstruct transactions and identify the beneficial owner. DNFBPs must ensure that records are available to regulators and law enforcement upon request. Failure to maintain records is a violation that can result in enforcement actions.

\vspace{0.1cm}
\noindent\textbf{Type of violation or warning signs:} A DNFBP discards records after one year. A DNFBP cannot produce records for a client when requested by a regulator. A DNFBP has no systematic approach to recordkeeping. A DNFBP's records are incomplete or inaccurate. A DNFBP does not retain copies of identification documents. A DNFBP does not have a records retention policy. A DNFBP destroys records before the required retention period.

\vspace{0.1cm}
\noindent\textbf{Correct answer approach:} The correct answer will state that DNFBPs must maintain records for at least five years and make them available to authorities. Eliminate answers that suggest shorter retention periods are acceptable or that records are optional. Remember: records are essential for law enforcement investigations. Without records, money laundering cannot be traced.
\end{frame}

% ============================================================
% SLIDE 12 - DNFBP AML PROGRAM REQUIREMENTS
% ============================================================
\begin{frame}{12. DNFBP AML Program Requirements}
\tiny
\noindent\textbf{Introduction to the topic:} Under FATF Recommendation 22, DNFBPs must establish and maintain an AML program that includes: written policies and procedures, a designated compliance officer, employee training, and independent testing (audit). The program must be commensurate with the DNFBP's size and risk profile. DNFBPs must also have procedures for identifying and reporting suspicious activity. The compliance officer must have access to all relevant information and the authority to act.

\vspace{0.1cm}
\noindent\textbf{Type of violation or warning signs:} A DNFBP has no written AML policies or procedures. A DNFBP has a compliance officer but that person has no authority. A DNFBP does not train employees on AML obligations. A DNFBP has never conducted an independent audit of its AML program. A DNFBP has no procedure for escalating suspicious activity to the compliance officer. A DNFBP's AML program does not address the specific risks of its business.

\vspace{0.1cm}
\noindent\textbf{Correct answer approach:} The correct answer will state that DNFBPs must have a written AML program with policies, a compliance officer, training, and independent testing. Eliminate answers that suggest smaller DNFBPs are exempt or that a program is optional. Remember: the AML program must be tailored to the DNFBP's specific risks and activities.
\end{frame}

% ============================================================
% SLIDE 13 - DNFBP SUSPICIOUS ACTIVITY INDICATORS
% ============================================================
\begin{frame}{13. DNFBP Suspicious Activity Indicators}
\tiny
\noindent\textbf{Introduction to the topic:} DNFBPs must be able to identify suspicious activity indicators specific to their sector. For real estate, indicators include cash purchases, rapid resale, and shell company buyers. For lawyers, indicators include complex structures with no legitimate purpose, use of client accounts to move funds, and refusals to disclose beneficial owners. For casinos, indicators include cash purchases just below reporting thresholds, "gambling" with small losses followed by cash-outs, and patrons from high-risk jurisdictions. Training employees on these indicators is essential.

\vspace{0.1cm}
\noindent\textbf{Type of violation or warning signs:} A real estate agent does not question a client's cash purchase of a \$10 million property. A lawyer does not question a client's request to create a complex trust with anonymous beneficiaries. A casino does not question a patron who gambles \$100 and cashes out \$10,000. A DNFBP employee sees obvious indicators but does not escalate them. A DNFBP does not provide training on suspicious activity indicators. A DNFBP does not have a list of red flags for its sector.

\vspace{0.1cm}
\noindent\textbf{Correct answer approach:} The correct answer will state that DNFBPs must train employees to identify sector-specific suspicious activity indicators and have procedures to escalate them. Eliminate answers that suggest employees should identify all suspicious activity without training or that training is optional. Remember: suspicious activity indicators are sector-specific — what is suspicious in real estate may not be suspicious in law.
\end{frame}

% ============================================================
% SLIDE 14 - DNFBP TIPPING OFF PROHIBITIONS
% ============================================================
\begin{frame}{14. DNFBP Tipping Off Prohibitions}
\tiny
\noindent\textbf{Introduction to the topic:} DNFBPs are prohibited from "tipping off" — informing a client that they are the subject of a SAR or that a SAR has been filed. Tipping off allows criminals to move assets, destroy evidence, or flee before law enforcement can act. The prohibition applies to DNFBPs as well as financial institutions. Tipping off is a criminal offense in most jurisdictions. DNFBPs must ensure that employees understand that they cannot discuss SAR filings with clients, even if the client asks directly.

\vspace{0.1cm}
\noindent\textbf{Type of violation or warning signs:} A lawyer tells a client that they have filed a SAR about the client's transaction. A real estate agent informs a buyer that the seller's suspicious activity has been reported. A casino employee tells a patron that they are "under observation." A DNFBP employee discusses SAR filing with colleagues who are not authorized to know. A DNFBP employee posts about a SAR filing on social media or in a public forum. A DNFBP employee denies the existence of a SAR when asked but does so in a way that implies its existence.

\vspace{0.1cm}
\noindent\textbf{Correct answer approach:} The correct answer will state that DNFBPs must not disclose that a SAR has been filed or that the client is under investigation. Eliminate answers that suggest DNFBPs can inform clients about SARs or that tipping off is only a violation for financial institutions. Remember: tipping off is a criminal offense. The client must not learn they are under suspicion until law enforcement acts.
\end{frame}

% ============================================================
% SLIDE 15 - SUPERVISION AND ENFORCEMENT OF DNFBPs
% ============================================================
\begin{frame}{15. Supervision and Enforcement of DNFBPs}
\tiny
\noindent\textbf{Introduction to the topic:} DNFBPs are subject to supervision and enforcement by designated authorities. In many jurisdictions, self-regulatory bodies (such as bar associations for lawyers, real estate boards, or casino regulators) conduct AML oversight. Regulators examine DNFBPs for compliance with AML obligations, including CDD, recordkeeping, SAR filing, and AML programs. Enforcement actions can include fines, license suspension, and criminal prosecution. DNFBPs that fail to comply with AML obligations face significant penalties.

\vspace{0.1cm}
\noindent\textbf{Type of violation or warning signs:} A DNFBP has never been examined or audited by a regulator. A DNFBP has no record of regulatory correspondence. A DNFBP has no AML program or an inadequate program. A DNFBP has never filed a SAR despite handling high-risk clients. A DNFBP does not respond to regulator inquiries. A DNFBP has been cited for AML deficiencies in a previous examination but has not corrected them. A DNFBP operates without any regulatory oversight.

\vspace{0.1cm}
\noindent\textbf{Correct answer approach:} The correct answer will state that DNFBPs are subject to supervision and enforcement and must comply with AML obligations. Eliminate answers that suggest DNFBPs are not supervised or that enforcement only applies to financial institutions. Remember: DNFBPs that fail to comply with AML obligations face fines, license suspension, and criminal liability.
\end{frame}

% ============================================================
% SUMMARY SLIDE - DNFBP AML PRINCIPLES
% ============================================================
\begin{frame}{Summary: Key Principles for DNFBP AML}
\tiny
\noindent\textbf{How to answer exam questions on DNFBPs:}

\vspace{0.1cm}
\noindent• \textbf{Coverage:} DNFBPs include lawyers, accountants, real estate agents, casinos, DPMs, and trust/company service providers.

\vspace{0.1cm}
\noindent• \textbf{CDD:} DNFBPs must identify clients and beneficial owners. No exemptions for "professional privilege" in financial transactions.

\vspace{0.1cm}
\noindent• \textbf{SAR Filing:} DNFBPs must file SARs when they suspect money laundering. Same standard as financial institutions.

\vspace{0.1cm}
\noindent• \textbf{Recordkeeping:} Records must be kept for at least five years. Essential for investigations.

\vspace{0.1cm}
\noindent• \textbf{AML Program:} Written policies, compliance officer, training, and independent testing required.

\vspace{0.1cm}
\noindent• \textbf{Tipping Off:} Prohibited — DNFBPs cannot disclose that a SAR has been filed.

\vspace{0.1cm}
\noindent• \textbf{Real Estate:} Cash purchases, shell company buyers, and rapid resale at loss = red flags.

\vspace{0.1cm}
\noindent• \textbf{Casinos:} Cash just under thresholds, "gambling" with losses then cashing out = red flags.

\vspace{0.1cm}
\noindent• \textbf{Lawyers:} Complex structures with no legitimate purpose, client accounts moving funds = red flags.

\vspace{0.1cm}
\noindent• \textbf{Trust Companies:} Anonymous beneficiaries, nominee arrangements without UBO identification = red flags.

\vspace{0.1cm}
\noindent• \textbf{Supervision:} DNFBPs are supervised and face enforcement for non-compliance.

\vspace{0.1cm}
\noindent• \textbf{Legal Privilege:} Does not exempt lawyers from AML obligations in financial transactions.

\vspace{0.3cm}
\centering
{Remember: DNFBPs are gatekeepers. They must know their clients, report suspicious activity, and cooperate with regulators.}
\end{frame}

\end{document}